\section{সংখ্যা পদ্ধতির বিকাশ}
গণনার শুরুতে প্রতিটি বস্তুর জন্য একটি দাগ বা পাথর রাখলেই কাজ চলত। কিন্তু বড় পরিমাণ, ভগ্নাংশ ও হিসাবের ফল সংরক্ষণ করতে সংক্ষিপ্ত প্রতীক এবং নির্দিষ্ট নিয়মের প্রয়োজন হয়। বিভিন্ন সভ্যতায় এই প্রয়োজন থেকেই সংখ্যা লেখার ভিন্ন পদ্ধতি গড়ে ওঠে।

\subsection{প্রতীক থেকে স্থানমান}
একটি প্রতীকের নির্দিষ্ট মান ধরে প্রয়োজনমতো সেটি পুনরাবৃত্ত করা যায়। আবার অল্প কয়েকটি অঙ্ককে ভিন্ন স্থানে বসিয়েও বড় সংখ্যা লেখা যায়। প্রথম ক্ষেত্রে প্রতীকের মান মুখ্য; দ্বিতীয় ক্ষেত্রে অঙ্কের অবস্থানও মান নির্ধারণ করে। আধুনিক গণনায় দ্বিতীয় ধারণাটি বিশেষ গুরুত্বপূর্ণ।
\begin{center}\begin{tabularx}{\linewidth}{|L{.18\linewidth}|Y|Y|}\hline
\rowcolor{ICTLightBlue}সভ্যতা/পদ্ধতি&সংখ্যা প্রকাশের বৈশিষ্ট্য&যা বোঝা প্রয়োজন\\\hline
মিশরীয়&১, ১০, ১০০ ইত্যাদির আলাদা প্রতীক; প্রয়োজনমতো প্রতীক যোগ&দশের গুচ্ছ ব্যবহার হলেও আধুনিক স্থানমান পদ্ধতি নয়\\\hline
ব্যাবিলনীয়&ষাটভিত্তিক স্থানমানের ব্যবহার&সময় ও কোণের একক বিভাজনে ৬০-এর ঐতিহাসিক প্রভাব আছে\\\hline
গ্রিক&বর্ণমালার অক্ষরের নির্দিষ্ট সংখ্যামান&সংখ্যা বোঝাতেও অক্ষর প্রতীক হতে পারে\\\hline
রোমান&I, V, X, L, C, D, M&দশের ঘাতভিত্তিক স্থানমান নেই; যোগ ও বিশেষ বিয়োগের নিয়ম আছে\\\hline
মায়ান&সাধারণ গণনায় ২০ ভিত্তি; বিন্দু, দাগ ও শূন্যের প্রতীক&দাগের মান ৫; সেটি পদ্ধতির ভিত্তি নয়\\\hline
ভারতীয় ও হিন্দু–আরবীয় ধারা&শূন্যসহ দশটি অঙ্ক এবং দশভিত্তিক স্থানমান&অল্প প্রতীকে বড় সংখ্যা লেখা ও নিয়মিত গণনা সহজ হয়\\\hline
\end{tabularx}\end{center}
ব্যাবিলনীয় ধারার সঙ্গে সম্পর্কিত পরিচিত বিভাজন: ১ ঘণ্টায় ৬০ মিনিট, ১ মিনিটে ৬০ সেকেন্ড; কৌণিক পরিমাপে ১ ডিগ্রিতে ৬০ মিনিট। মায়ান সাধারণ গণনার স্থানমান ১, ২০, ৪০০ ইত্যাদি; ক্যালেন্ডারের কিছু গণনায় বিশেষ স্থানমান ব্যবহৃত হয়।

\smallhead{উদাহরণ: রোমান প্রতীকের মান ও বিন্যাস}
\begin{center}\begin{tabular}{|c|c|c|c|c|c|c|c|}\hline
প্রতীক&I&V&X&L&C&D&M\\\hline
মান&1&5&10&50&100&500&1000\\\hline
\end{tabular}\end{center}
XVII-তে $10+5+1+1=17$। IX-এ I আগে থাকায় বিশেষ নিয়মে $10-1=9$; XI-এ $10+1=11$। সুতরাং রোমান পদ্ধতিতে দশমিকের মতো স্থানমান নেই, কিন্তু প্রতীকগুলোর ক্রম একেবারেই গুরুত্বহীন—এ কথা বলা ঠিক নয়। ঘড়ির ডায়াল ও অধ্যায়ের ক্রম প্রকাশে এখনো রোমান সংখ্যা দেখা যায়।

\subsection{শূন্য ও আধুনিক স্থানমান}
শূন্য একটি সংখ্যামান; একই সঙ্গে এটি কোনো স্থানে পরিমাণ না থাকার চিহ্ন হিসেবেও কাজ করে। যেমন
\[407=4\times100+0\times10+7\times1.\]
এখানে শূন্যটি দশকের স্থান ধরে রেখেছে। এটি বাদ দিয়ে 47 লিখলে ৪-এর অবদান ৪০০ থেকে ৪০ হয়। আবার $4.07$-এ ভগ্নাংশের শুরুর শূন্য বাদ দিয়ে $4.7$ লিখলেও মান বদলায়।

শূন্যসহ দশভিত্তিক স্থানমানের বিকাশে ভারতীয় গণিতচর্চার গুরুত্বপূর্ণ অবদান ছিল; আরব বিশ্বের গণিতচর্চা ও লেখালেখির মাধ্যমে পদ্ধতিটি আরও বিস্তার লাভ করে। তাই আধুনিক দশমিক পদ্ধতিকে হিন্দু–আরবীয় পদ্ধতিও বলা হয়। কম্পিউটার ব্যবস্থায় একই স্থানমানের ধারণা ভিত্তি ২, ৮ ও ১৬-তেও প্রয়োগ করা হয়।

\smallhead{ধারণা যাচাই}
\textbf{প্রশ্ন:} মায়ান প্রতীকে একটি দাগের মান ৫ হলে পদ্ধতির ভিত্তি কি ৫?\\
\textbf{উত্তর:} না। প্রতীকের নিজস্ব মান ও পদ্ধতির ভিত্তি পৃথক বিষয়। সাধারণ মায়ান গণনায় ভিত্তি ২০; বিন্দুর মান ১ এবং দাগের মান ৫।
